\begin{theorem}[Extremal envelope contact]
Let $F$ be a finite-dimensional real inner-product space, write $m$ for its
dimension and let $1\le C_0<\infty$. Suppose that $\Aclass_{m,C_0}$, in the
sense of \noderef{IsFitClass}, contains a member $v^{(0)}$ that is not affine on
space-time in the sense of \noderef{IsSpaceTimeAffine}. Then there are
$\kappa\in(0,\infty)$, a real number $H$ and a function $v\in\Aclass_{m,C_0}$
centered at $z_*=(0,1)$ with constant $C_0$ in the sense of
\noderef{IsCenteredAt}, such that $v$ satisfies the horizontal Jensen bound with
constant $\kappa$ at every time, in the sense of \noderef{HasJensenBound} with
$I=\R$, and such that the contact envelope
\[
  W(x,t)=\Cenv_tv(x,t)-\frac{\kappa^2t}{2}\qquad(t>0)
\]
of \noderef{ContactEnv}, formed from the barycentric envelope of
\noderef{BaryEnv}, has the following properties.
\begin{enumerate}
\item $W$ is continuous on $\{(x,t):t>0\}$ and has subquadratic spatial growth
on compact subintervals of $(0,\infty)$, in the sense of
\noderef{SubquadraticGrowth}.
\item $W$ is a viscosity subsolution of the limiting equation on
$\{(x,t):t>0\}$, in the sense of \noderef{IsLimitSubOn}.
\item For every $t>0$,
\[
  v(x,t)-\frac{\kappa^2t}{2}\le W(x,t)\le v(x,t),
\]
and $W(z_*)=v(z_*)=0$.
\item For every $t>0$ the function $x\mapsto W(x,t)-|x|^2/(2t)$ is concave on
$F$.
\item At every contact with $t>0$ the barycentric gap is exactly
$\kappa^2t/2$: if $W(x,t)=v(x,t)$ and $t>0$ then
\[
  \frac{\Cenv_tv(x,t)-v(x,t)}{t}=\frac{\kappa^2}{2}.
\]
\item The envelope Hamiltonian of \noderef{EnvHamiltonian} is finite at $z_*$,
with value $H_v(z_*)=H$, and the forward envelope slopes of
\noderef{EnvSlopeP} satisfy
\[
  \lim_{\ell\downarrow0}\Senv_\ell^+W(z_*)=H .
\]
\end{enumerate}
\end{theorem}
\begin{proof}
Throughout, $z_*=(0,1)$, a \emph{configuration} means a finite horizontal
probability configuration in the sense of \noderef{IsBaryConfig}, and its
\emph{radius} is the root mean square radius of \noderef{BaryRadius}. For a
function $g$ in $\Aclass_{m,C_0}$ we use freely that $g$ is continuous on all of
space-time and is a viscosity solution of the limiting equation
$F[g]=D^2g[(\nabla_xg,1)]=0$ there, by \noderef{IsFitClass} and
\noderef{IsLimitSolutionOn}. By \noderef{UniformSubquadratic}, applied with
$z_c=(0,0)$ and with $z_c=(0,1)$, every member of $\Aclass_{m,C_0}$ centered at
one of those two points has subquadratic spatial growth on compact time
intervals uniformly over the class; in particular each such member satisfies
\noderef{SubquadraticGrowth} with $I=\R$, hence also with $I=(0,\infty)$, since
decreasing $I$ only decreases the family of compact subintervals tested.

\emph{The envelope bound pairs.} Whenever a real function $f$ on space-time is
continuous on $\Omega_I=\{(x,t):t\in I\}$ and has subquadratic spatial growth on
compact subintervals of $I$ in the sense of \noderef{SubquadraticGrowth}, we
apply \noderef{EnvelopeCompetitorBounds} to $f$ with that $I$. Its three
conjuncts are exactly the three facts used below, and we name them in that
order. For $s>0$ and $z=(x,t)$ with $t\in I$ write
\[
  V(f,s,z)=\Bigl\{\sum_i\theta_if(y_i,t)-\frac{\rho^2}{2s}:
     (\theta_i,y_i)\text{ a configuration with mean }x\text{ and radius }
     \rho\Bigr\}
\]
for the set of barycentric competitor values at $z$, the configurations being
those of \noderef{IsBaryConfig} and the radius that of \noderef{BaryRadius}.
\begin{enumerate}
\item[(B1)] The third conjunct of \noderef{EnvelopeCompetitorBounds}: every
element of $V(f,s,z)$ is at most $\Cenv_sf(z)$, and $\Cenv_sf(z)$ is at most
every real number that dominates every element of $V(f,s,z)$.
\item[(B2)] Its first conjunct: for $\ell>0$ with $t\in I$ and $t+\ell\in I$,
every $f(y,t+\ell)-|y-x|^2/(2\ell)$ with $y\in F$ is at most $\Penv_\ell f(z)$,
and $\Penv_\ell f(z)$ is at most every real number dominating all of them.
\item[(B3)] Its second conjunct: for $\ell>0$ with $t\in I$ and $t-\ell\in I$,
$\Menv_\ell f(z)$ is at most every $f(y,t-\ell)+|y-x|^2/(2\ell)$ with $y\in F$,
and every real number below all of them is at most $\Menv_\ell f(z)$.
\end{enumerate}
Two remarks on how these are used. First, the domination half of each pair
exhibits the envelope as an upper bound of its competitor set in (B1) and (B2)
and as a lower bound in (B3); so each of those sets is bounded on the relevant
side by a real number. Nonemptiness is a separate and immediate matter: the
one-point configuration $N=1$, $\theta_1=1$, $y_1=x$ has mean $x$ and radius
$0$ and is admissible by \noderef{IsBaryConfig}, so $f(z)\in V(f,s,z)$, and the
choice $y=x$ is always available in (B2) and (B3). Second, no separate
finiteness assertion is needed anywhere below: $\Cenv_s$, $\Penv_\ell$ and
$\Menv_\ell$ are real-valued by \noderef{BaryEnv}, \noderef{QuadEnvP} and
\noderef{QuadEnvM} together with the preamble's convention for the supremum or
infimum of an empty or unbounded set, and what that convention withholds --
namely that the value is the genuine least upper bound, respectively greatest
lower bound, of its competitor set rather than the conventional substitute --
is exactly what the two halves of each pair supply.
\emph{Changes of frame act on configurations.} Let $g$ be a real function on
space-time, let $z_0=(x_0,t_0)$, $R>0$, $p_0\in F$, $a_0\in\R$, and let $\bar g$
be the change of frame of \noderef{FrameChange} based at $z_0$ with scale $R$
and parameters $(p_0,a_0)$,
\[
  \bar g(y,s)=\frac{g\bigl(x_0+R(y+p_0s),\,t_0+Rs\bigr)-g(z_0)}{R}
     -p_0\cdot y-\Bigl(a_0+\frac{|p_0|^2}{2}\Bigr)s .
\]
Fix $s\in\R$ and put $T_s(y)=x_0+R(y+p_0s)$ and $t=t_0+Rs$, so that
$\bar g(y,s)=[g(T_s(y),t)-g(z_0)]/R-p_0\cdot y-(a_0+|p_0|^2/2)s$. We write
(F0) for \noderef{FrameConfigCovariance} applied to these data. Its parts (i),
(ii) and (iii) say that $T_s$ carries a family of weights and positions with
mean $\bar y$ onto one with mean $T_s(\bar y)$ and back, in the sense of
\noderef{IsBaryConfig}, keeping the weights, multiplying the radius of
\noderef{BaryRadius} by $R$, and hitting every point of $F$; and its part (iv)
says that for every configuration $(\theta_i,y_i)$ with mean $\bar y$,
\[
  \sum_i\theta_i\bar g(y_i,s)-\bar g(\bar y,s)
    =\frac1R\Bigl(\sum_i\theta_ig(T_s(y_i),t)-g(T_s(\bar y),t)\Bigr).
\]

Three consequences will be used.
\begin{enumerate}
\item[(F1)] If $g$ satisfies the Jensen bound with constant $\kappa\ge0$ at
every time (\noderef{HasJensenBound} with $I=\R$), then so does $\bar g$: by
(F0)(iv), together with the radius identity (F0)(ii), for a configuration with
mean $y$ and radius $\rho$,
\[
  \sum_i\theta_i\bar g(y_i,s)-\bar g(y,s)
   =\frac1R\Bigl(\sum_i\theta_ig(T_sy_i,t)-g(T_sy,t)\Bigr)
   \le\frac{\kappa\,R\rho}{R}=\kappa\rho .
\]
\item[(F2)] Suppose that $g$ and $\bar g$ are continuous on space-time with
subquadratic spatial growth on compact time intervals. Then for every $\sigma>0$
and every $(y,s)$, with $Z=(T_s(y),t)$,
\[
  \Cenv_\sigma\bar g(y,s)-\bar g(y,s)
    =\frac{\Cenv_{R\sigma}g(Z)-g(Z)}{R}.
\]
Indeed, by (F0)(iv), the element of $V(\bar g,\sigma,(y,s))$ attached to a
configuration with mean $y$ and radius $\rho$, minus $\bar g(y,s)$, equals
\[
  \frac1R\Bigl(\sum_i\theta_ig(T_sy_i,t)-g(Z)\Bigr)-\frac{\rho^2}{2\sigma}
   =\frac1R\Bigl(\sum_i\theta_ig(T_sy_i,t)-g(Z)
       -\frac{(R\rho)^2}{2R\sigma}\Bigr),
\]
because $(R\rho)^2/(2R^2\sigma)=\rho^2/(2\sigma)$; and by (F0)(i),(ii),(iii) the
map
$(\theta_\bullet,y_\bullet)\mapsto(\theta_\bullet,T_sy_\bullet)$ is a bijection
of the configurations occurring in $V(\bar g,\sigma,(y,s))$ onto those occurring
in $V(g,R\sigma,Z)$, carrying radius $\rho$ to $R\rho$. Thus
$V(\bar g,\sigma,(y,s))-\bar g(y,s)$ is the image of
$V(g,R\sigma,Z)-g(Z)$ under multiplication by the positive number $1/R$. Both
sets are nonempty, as noted above, and the domination half of (B1) exhibits
each of $\Cenv_\sigma\bar g(y,s)$ and $\Cenv_{R\sigma}g(Z)$ as an upper bound of
its own set while the extremality half makes it the least such bound; a positive
dilation of the reals carries the least upper bound of a nonempty set bounded
above to the least upper bound of the dilated set, so the two suprema correspond
in the same way, which is (F2).

Both applications of (B1) are legitimate: the hypotheses of
\noderef{EnvelopeCompetitorBounds} hold for $g$ with the ambient $I$ and for
$\bar g$ with the corresponding one, by the continuity and growth assumptions in
force wherever (F2) is invoked.
\item[(F3)] If $\bar g$ is affine on space-time in the sense of
\noderef{IsSpaceTimeAffine} then so is $g$. Indeed the definition of $\bar g$
rearranges to
\[
  g\bigl(x_0+R(y+p_0s),t_0+Rs\bigr)
   =g(z_0)+R\Bigl[\bar g(y,s)+p_0\cdot y
        +\Bigl(a_0+\frac{|p_0|^2}{2}\Bigr)s\Bigr],
\]
whose right side is an affine function of $(y,s)$ when $\bar g$ is, while
$(y,s)\mapsto(x_0+R(y+p_0s),t_0+Rs)$ is an affine bijection of space-time, with
inverse $(x,t)\mapsto\bigl((x-x_0)/R-p_0(t-t_0)/R,\ (t-t_0)/R\bigr)$. Composing
with that inverse exhibits $g$ as an affine function.
\end{enumerate}

\emph{Claim T: translation in time.}
Let $c\in\R$, let $g$ be a real function on space-time and put
$g_c(x,t)=g(x,t-c)$. Then:
\begin{enumerate}
\item[(i)] $g_c$ is continuous iff $g$ is, and $g_c$ is a viscosity solution of
the limiting equation on all of space-time iff $g$ is
(\noderef{IsLimitSolutionOn});
\item[(ii)] $(p,a)$ is a fit for $g_c$ at $(x,t)$ of radius $\rho$ with constant
$C_0$ iff it is a fit for $g$ at $(x,t-c)$ of radius $\rho$ with constant $C_0$
(\noderef{IsFit});
\item[(iii)] $g_c$ satisfies the Jensen bound with constant $\kappa$ at every
time iff $g$ does (\noderef{HasJensenBound} with $I=\R$);
\item[(iv)] $\Cenv_sg_c(x,t)=\Cenv_sg(x,t-c)$ for every $s>0$ and $(x,t)$;
\item[(v)] consequently $g_c\in\Aclass_{m,C_0}$ iff $g\in\Aclass_{m,C_0}$, and
$g_c$ is centered at $(x,t)$ iff $g$ is centered at $(x,t-c)$.
\end{enumerate}
\emph{Proof.} (ii) holds because \noderef{IsFit} involves only the base point
and increments $(h,\ell)$ from it, and the base point of $g_c$ at $(x,t)$ is the
base point of $g$ at $(x,t-c)$, with identical increments. (iii) and (iv) hold
because \noderef{HasJensenBound} and \noderef{BaryEnv} involve, for each
assertion, only one horizontal slice, and the slice of $g_c$ at time $t$ is the
slice of $g$ at time $t-c$; in (iii) the time quantifier ranges over all of $\R$
in both cases. For (i), write $\tau(x,t)=(x,t-c)$, so $g_c=g\circ\tau$ and
$\tau$ is a bijective isometry of space-time whose inverse is the translation
$(x,t)\mapsto(x,t+c)$; continuity is clear. If $\phi$ is twice continuously
differentiable at $\hat z$ and $g_c-\phi$ has a local maximum at $\hat z$, then
$g-\phi\circ\tau^{-1}$ has a local maximum at $\tau(\hat z)$ and
$\phi\circ\tau^{-1}$ is twice continuously differentiable at $\tau(\hat z)$,
with the same gradient and second derivative as $\phi$ at $\hat z$, because
$\tau^{-1}$ is a translation. Hence $F[\phi\circ\tau^{-1}](\tau\hat z)
=F[\phi](\hat z)$ by \noderef{LimitOp}, and the subsolution inequality of
\noderef{IsLimitSubOn} transfers; minima and \noderef{IsLimitSuperOn} are
identical, and the converse direction follows by replacing $c$ with $-c$. Part
(v) combines (i), (ii) with the definition of \noderef{IsFitClass} and
\noderef{IsCenteredAt}. This proves Claim T.

\emph{Step 1: a centered nonaffine member and the extremal constant.}
By \noderef{IsFitClass} choose a fit $(p_0,a_0)$ for $v^{(0)}$ at the origin of
radius $1$, and let $w$ be the change of frame of \noderef{FrameChange} of
$v^{(0)}$ based at $z_0=0$ with scale $R=1$ and parameters $(p_0,a_0)$. By the
first clause of \noderef{FrameInvariance}, $w\in\Aclass_{m,C_0}$. Its second
clause, with $\zeta=0$ and $\rho=1$, has as hypothesis exactly a fit for
$v^{(0)}$ at the origin of radius $1\cdot1$, namely $(p_0,a_0)$, and yields that
$(p_0-p_0,a_0-a_0)=(0,0)$ is a fit for $w$ at the origin of radius $1$; and
$w(0)=0$ directly from \noderef{FrameChange}. So $w$ is centered at the origin
with constant $C_0$ in the sense of \noderef{IsCenteredAt}. Moreover $w$ is not
affine: otherwise $v^{(0)}$ would be affine by (F3).

Now \noderef{JensenSupremumFinite}, applied to $C_0$ and $w$, supplies a
function $\check v\in\Aclass_{m,C_0}$, centered at the origin and not affine,
together with a constant $\kappa$ such that
\begin{equation}\label{eq:ec-kappa}
  0<\kappa<\infty,\qquad
  \check v\text{ satisfies the Jensen bound with constant }\kappa
  \text{ at every time},
\end{equation}
and such that $\kappa\le\kappa'$ for every $\kappa'$ with which $\check v$
satisfies that bound. This $\kappa$ is the extremal constant of the proposition.

\emph{Step 2: near-attaining configurations.}
For every $\varepsilon>0$ there are a time $t$ and a configuration
$(\theta_i,y_i)$ with mean $x$ and radius $\rho>0$ such that
\[
  \sum_i\theta_i\check v(y_i,t)-\check v(x,t)>(\kappa-\varepsilon)\rho .
\]
Otherwise $\check v$ would satisfy the Jensen bound with the constant
$\kappa-\varepsilon$: for configurations of positive radius this is the failure
of the displayed inequality, and for a configuration of radius $0$ all $y_i$
equal $x$, because $\sum_i\theta_i|y_i-x|^2=0$ with every $\theta_i>0$, so the
left side vanishes and the required inequality reads $0\le0$. The minimality in
Step 1, applied to the admissible constant $\kappa'=\kappa-\varepsilon$, would
then give $\kappa\le\kappa-\varepsilon$, which is false for $\varepsilon>0$.
Choose accordingly, for each integer $j\ge0$, a time $t_j$ and a configuration
$(\theta_{j,i},y_{j,i})_{i\le N_j}$ with mean $x_j$ and radius $\rho_j>0$ such
that
\begin{equation}\label{eq:ec-near}
  D_j:=\sum_i\theta_{j,i}\check v(y_{j,i},t_j)-\check v(x_j,t_j)
     >\Bigl(\kappa-\frac1{j+1}\Bigr)\rho_j .
\end{equation}

\emph{Step 3: normalizing each configuration.}
Fix $j$. By \noderef{IsFitClass} choose a fit $(p_j,a_j)$ for $\check v$ at
$(x_j,t_j)$ of radius $\rho_j>0$, and let $v_j$ be the change of frame
\noderef{FrameChange} of $\check v$ based at $z_j=(x_j,t_j)$ with scale
$R=\rho_j$ and parameters $(p_j,a_j)$. As in Step 1,
\noderef{FrameInvariance} gives $v_j\in\Aclass_{m,C_0}$ and, with $\zeta=0$ and
$\rho=1$, that $(0,0)$ is a fit for $v_j$ at the origin of radius $1$, while
$v_j(0)=0$; so $v_j$ is centered at the origin. By (F1) and
\eqref{eq:ec-kappa}, $v_j$ satisfies the Jensen bound with the same constant
$\kappa$ at every time.

Put $\hat y_{j,i}=(y_{j,i}-x_j)/\rho_j$. In (F0) for this frame change, take
$s=0$; then $T_0(y)=x_j+\rho_jy$ and $t=t_j$, so $T_0(\hat y_{j,i})=y_{j,i}$ and
$T_0(0)=x_j$. By (F0)(i),(ii), $(\theta_{j,i},\hat y_{j,i})$ is a configuration
with mean $0$ and radius $1$, since its image under $T_0$ has mean $x_j$ and
radius $\rho_j=\rho_j\cdot1$. By (F0)(iv) and $v_j(0,0)=0$,
\[
  \sum_i\theta_{j,i}v_j(\hat y_{j,i},0)
    =\sum_i\theta_{j,i}v_j(\hat y_{j,i},0)-v_j(0,0)
    =\frac{D_j}{\rho_j}>\kappa-\frac1{j+1},
\]
while the Jensen bound for $v_j$ at time $0$ with this configuration gives
$\sum_i\theta_{j,i}v_j(\hat y_{j,i},0)\le\kappa\cdot1$. Hence
\begin{equation}\label{eq:ec-tendsto}
  \sum_i\theta_{j,i}v_j(\hat y_{j,i},0)\longrightarrow\kappa
  \qquad(j\to\infty).
\end{equation}

\emph{Step 4: compactness and the attainment of the sandwich.}
By \noderef{ClassCompactness}, applied with $z_c=(0,0)$ to the sequence
$(v_j)$, there are a strictly increasing $\iota$ and a function
$\hat v\in\Aclass_{m,C_0}$, centered at the origin, with $v_{\iota(j)}\to\hat v$
uniformly on every compact set. For each fixed time and each fixed configuration
the Jensen inequality for $v_{\iota(j)}$ passes to the limit of the finitely
many values involved, so $\hat v$ satisfies the Jensen bound with constant
$\kappa$ at every time (\noderef{HasJensenBound} with $I=\R$).

\emph{Claim 4a.} $\displaystyle
  \sum_i\theta_{\iota(j),i}\hat v(\hat y_{\iota(j),i},0)\to\kappa$.
\emph{Proof.} Let $\delta>0$ and put $\varepsilon=\delta/8$. By
\noderef{UniformSubquadratic} with $z_c=(0,0)$, the compact time set $J=\{0\}$
and this $\varepsilon$, there is $M>0$ such that $|g(y,0)|\le\varepsilon|y|^2$
whenever $|y|\ge M$, for every $g\in\Aclass_{m,C_0}$ centered at the origin;
this applies to every $v_j$ and to $\hat v$. Splitting the weighted sum over the
indices with $|\hat y_{j,i}|<M$ and those with $|\hat y_{j,i}|\ge M$, and using
$\sum_i\theta_{j,i}|\hat y_{j,i}|^2=1$ (the configuration has mean $0$ and
radius $1$),
\[
  \Bigl|\sum_i\theta_{j,i}\bigl(\hat v-v_j\bigr)(\hat y_{j,i},0)\Bigr|
   \le\sup_{|y|\le M}\bigl|\hat v(y,0)-v_j(y,0)\bigr|
      +2\varepsilon\sum_i\theta_{j,i}|\hat y_{j,i}|^2
   \le\sup_{|y|\le M}\bigl|\hat v(y,0)-v_j(y,0)\bigr|+2\varepsilon .
\]
The set $\{(y,0):|y|\le M\}$ is compact, so the first term tends to $0$ along
$\iota$, and for large $j$ the right side is at most
$\delta/2+\delta/4<\delta$. Together with \eqref{eq:ec-tendsto} this proves
Claim 4a.

Put $R=1/\kappa>0$. The configuration
$(\theta_{\iota(j),i},\hat y_{\iota(j),i})$ has mean $0$ and radius $1$, so it
occurs in the set $V(\hat v,R,(0,0))$, and the domination half of (B1) applies
to it, the hypotheses of \noderef{EnvelopeCompetitorBounds} holding for
$\hat v$ with $I=\R$ by the opening paragraph. Hence
\[
  \Cenv_R\hat v(0,0)\ \ge\
    \sum_i\theta_{\iota(j),i}\hat v(\hat y_{\iota(j),i},0)-\frac{1}{2R},
\]
and letting $j\to\infty$, using Claim 4a, $\hat v(0,0)=0$ and
$1/(2R)=\kappa/2$,
\[
  \Cenv_R\hat v(0,0)-\hat v(0,0)\ \ge\ \kappa-\frac\kappa2=\frac\kappa2 .
\]
In the opposite direction \noderef{BaryEnvSandwich}, applied with $I=\R$,
$f=\hat v$ and the constant $\kappa$, gives
$\Cenv_R\hat v(0,0)-\hat v(0,0)\le\kappa^2R/2=\kappa/2$. Therefore
\begin{equation}\label{eq:ec-attain}
  \Cenv_{1/\kappa}\hat v(0,0)-\hat v(0,0)=\frac\kappa2 .
\end{equation}

\emph{Step 5: the final frame and the normalized contact.}
By \noderef{IsFitClass} choose a fit $(p_*,a_*)$ for $\hat v$ at the origin of
radius $R=1/\kappa$, let $g$ be the change of frame \noderef{FrameChange} of
$\hat v$ based at $z_0=0$ with scale $R$ and parameters $(p_*,a_*)$, and define
\[
  v(y,t)=g(y,t-1)
   =\frac{\hat v\bigl(R(y+p_*(t-1)),R(t-1)\bigr)-\hat v(0,0)}{R}
     -p_*\cdot y-\Bigl(a_*+\frac{|p_*|^2}{2}\Bigr)(t-1),
\]
the single change of frame and time translation of the proposition. By
\noderef{FrameInvariance}, $g\in\Aclass_{m,C_0}$, and its second clause with
$\zeta=0$, $\rho=1$ --- whose hypothesis is the chosen fit for $\hat v$ at the
origin of radius $R\cdot1$ --- shows that $(0,0)$ is a fit for $g$ at the origin
of radius $1$, while $g(0)=0$; so $g$ is centered at the origin. By (F1), $g$
satisfies the Jensen bound with constant $\kappa$ at every time. By Claim T with
$c=1$ it follows that
\begin{equation}\label{eq:ec-final}
  v\in\Aclass_{m,C_0},\quad v\text{ is centered at }z_*=(0,1),\quad
  v\text{ satisfies the Jensen bound with constant }\kappa
  \text{ at every time},
\end{equation}
and in particular $v(z_*)=0$ by \noderef{IsCenteredAt}.

Both $g$ and $\hat v$ are continuous with subquadratic spatial growth on compact
time intervals, by the opening paragraph applied to the centered members $g$ and
$\hat v$ of the class. So (F2) applies with $\sigma=1$ and $(y,s)=(0,0)$, for
which $T_0(0)=0$ and $t=0$, that is $Z=(0,0)$:
\[
  \Cenv_1g(0,0)-g(0,0)=\frac{\Cenv_R\hat v(0,0)-\hat v(0,0)}{R}
    =\kappa\cdot\frac{\kappa}{2}=\frac{\kappa^2}{2},
\]
using \eqref{eq:ec-attain} and $1/R=\kappa$. By Claim T(iv) with $c=1$,
$\Cenv_1v(0,1)=\Cenv_1g(0,0)$, and $v(0,1)=g(0,0)$; hence
\begin{equation}\label{eq:ec-contact}
  \Cenv_1v(0,1)-v(0,1)=\frac{\kappa^2}{2}.
\end{equation}

\emph{Step 6: the contact envelope.}
Let $W$ be the contact envelope of \noderef{ContactEnv} formed from $v$ and
$\kappa$, so $W(x,t)=\Cenv_tv(x,t)-\kappa^2t/2$ for $t>0$. By
\eqref{eq:ec-final} and the opening paragraph, $v$ is continuous on space-time,
satisfies \noderef{SubquadraticGrowth} with $I=\R$ and with $I=(0,\infty)$, and
is a viscosity solution of the limiting equation everywhere.

\emph{(a) The sandwich and the contact.} \noderef{BaryEnvSandwich}, applied with
$I=(0,\infty)$, $f=v$ and the constant $\kappa$, gives for every $s>0$ and every
$z=(x,t)$ with $t>0$ that $v(z)\le\Cenv_sv(z)\le v(z)+\kappa^2s/2$. Taking
$s=t$ and subtracting $\kappa^2t/2$,
\begin{equation}\label{eq:ec-sandwich}
  v(x,t)-\frac{\kappa^2t}{2}\le W(x,t)\le v(x,t)\qquad(t>0).
\end{equation}
At $z_*$, \eqref{eq:ec-contact} gives
$W(z_*)=\Cenv_1v(0,1)-\kappa^2/2=v(0,1)=0$, so $W(z_*)=v(z_*)=0$.

\emph{(b) Continuity and semiconcavity.} Apply \noderef{BaryEnvDuality} with
$I=(0,\infty)$ and $f=v$; its hypotheses are the continuity and growth just
recorded. Its continuity clause gives that
$(s,z)\mapsto\Cenv_sv(z)$ is continuous at every pair $(s,z_0)$ with $s>0$ and
the time of $z_0$ positive; composing with the continuous map
$z\mapsto(t(z),z)$, where $t(z)$ is the time coordinate, shows that
$z\mapsto\Cenv_{t(z)}v(z)$ is continuous on $\{t>0\}$, and subtracting the
continuous function $\kappa^2t/2$ gives the continuity of $W$ there. Its
concavity clause, applied with $s=t$ for a fixed $t>0$, gives that
$x\mapsto\Cenv_tv(x,t)-|x|^2/(2t)$ is concave on $F$; subtracting the constant
$\kappa^2t/2$, so is $x\mapsto W(x,t)-|x|^2/(2t)$, for every $t>0$.

\emph{(c) Growth.} Let $J\subseteq(0,\infty)$ be compact and let
$\varepsilon>0$. As $J$ is bounded there is $c_J$ with $\kappa^2t/2\le c_J$ for
$t\in J$. By \noderef{SubquadraticGrowth} for $v$ with $I=(0,\infty)$ and
$\varepsilon/2$ there is $M_1>0$ with $|v(x,t)|\le(\varepsilon/2)|x|^2$ whenever
$|x|\ge M_1$ and $t\in J$. By \eqref{eq:ec-sandwich},
$|W(x,t)|\le|v(x,t)|+\kappa^2t/2$ for $t>0$, so for $t\in J$ and
$|x|\ge\max\{M_1,\sqrt{2c_J/\varepsilon}\}$ we get
$|W(x,t)|\le(\varepsilon/2)|x|^2+c_J\le\varepsilon|x|^2$. This is
\noderef{SubquadraticGrowth} for $W$ with $I=(0,\infty)$.

\emph{(d) The subsolution property.} First, $v$ is a viscosity subsolution of
the limiting equation on the open set $\{t>0\}$ in the sense of
\noderef{IsLimitSubOn}: it is continuous there, and if $\phi$ is twice
continuously differentiable at a point $z$ with positive time and $v-\phi$ has a
local maximum at $z$ relative to $\{t>0\}$, then, because $\{t>0\}$ is open and
contains $z$, a neighbourhood of $z$ witnessing that relative maximum may be
intersected with $\{t>0\}$ and is then a neighbourhood of $z$ in the whole
space on which the inequality holds; so $v-\phi$ has a local maximum at $z$
relative to the whole space, and $0\le F[\phi](z)$ follows from
\noderef{IsLimitSubOn} for $v$ on the whole space. Now
\noderef{BarycentricSubsolution}, applied with $f=v$, with this subsolution
property, with \noderef{SubquadraticGrowth} on $I=(0,\infty)$, and with
$c=\kappa^2/2$, shows that
$z\mapsto\Cenv_{t(z)}v(z)-(\kappa^2/2)t(z)$, which is exactly $W$ by
\noderef{ContactEnv}, is a viscosity subsolution of the limiting equation on
$\{t>0\}$.

\emph{(e) The gap at a contact.} If $t>0$ and $W(x,t)=v(x,t)$ then
\noderef{ContactEnv} gives $\Cenv_tv(x,t)-v(x,t)=\kappa^2t/2$, hence
\[
  \frac{\Cenv_tv(x,t)-v(x,t)}{t}=\frac{\kappa^2}{2}>0 .
\]

\emph{Step 7: the slope at the contact.}
Since $v$ is continuous, has subquadratic spatial growth with $I=\R$ and solves
the limiting equation on all of space-time, \noderef{EnvHamiltonianSolution}
applies with $I=\R$ and $f=v$. Its third clause gives $H_v(z_*)\ne+\infty$,
where $H_v$ is the envelope Hamiltonian of \noderef{EnvHamiltonian}; its first
clause gives that $\ell\mapsto\Senv_\ell^+v(z_*)$ and
$\ell\mapsto\Senv_\ell^-v(z_*)$ are nondecreasing on $(0,\infty)$
(\noderef{EnvSlopeP}, \noderef{EnvSlopeM}); and its second clause, whose
hypothesis holds because $I=\R$, gives
\begin{equation}\label{eq:ec-ham}
  \inf_{\ell>0}\Senv_\ell^-v(z_*)=H_v(z_*)
   =\inf_{\ell>0}\Senv_\ell^+v(z_*),
\end{equation}
the second equality being the definition in \noderef{EnvHamiltonian}.
By \eqref{eq:ec-final} and \eqref{eq:ec-contact}, the hypotheses of
\noderef{ContactFiniteSlope} hold for $C_0$, $\kappa>0$ and $v$, so
$H_v(z_*)\ne-\infty$. Hence there is a real number $H$ with $H_v(z_*)=H$, and by
\eqref{eq:ec-ham} every $\Senv_\ell^\pm v(z_*)$ is a real number not smaller
than $H$.

Next we compare the slopes of $W$ and $v$ at $z_*$, where $W(z_*)=v(z_*)$ by
Step 6(a). By Step 6(b),(c) the function $W$ satisfies the hypotheses of
\noderef{EnvelopeCompetitorBounds} with $I=(0,\infty)$, and $v$ satisfies them
with $I=\R$; the relevant times $1\pm\ell$ lie in $(0,\infty)$ for $0<\ell<1$,
and $1+\ell\in(0,\infty)$ for every $\ell>0$. So (B2) applies to
$\Penv_\ell W(z_*)$ and $\Penv_\ell v(z_*)$ for every $\ell>0$, and (B3) to
$\Menv_\ell W(z_*)$ and $\Menv_\ell v(z_*)$ for $0<\ell<1$: each of the four is
the genuine least upper bound, respectively greatest lower bound, of its
competitor set. All four are real numbers, and hence so are
$\Senv_\ell^+W(z_*)$ and, for $0<\ell<1$, $\Senv_\ell^-W(z_*)$, by the second
remark of the opening paragraph.

For $\ell>0$, the set defining $\Penv_\ell$ at $z_*$ in \noderef{QuadEnvP}
involves only values at time $1+\ell>0$, where $W\le v$ by
\eqref{eq:ec-sandwich}; so each of its elements for $W$ is at most the element
for $v$ with the same $y$, and the latter set is bounded above by the
domination half of (B2) for $v$; the extremality half of (B2) for $W$ then gives
the comparison.
Hence $\Penv_\ell W(z_*)\le\Penv_\ell v(z_*)$ and, by \noderef{EnvSlopeP} and
$W(z_*)=v(z_*)$,
\begin{equation}\label{eq:ec-slope-plus}
  \Senv_\ell^+W(z_*)\le\Senv_\ell^+v(z_*)\qquad(\ell>0).
\end{equation}
For $0<\ell<1$, the set defining $\Menv_\ell$ at $z_*$ in \noderef{QuadEnvM}
involves only values at time $1-\ell>0$, where again $W\le v$; so
$\Menv_\ell W(z_*)$, a lower bound for the $W$-set, is a lower bound for the
$v$-set as well, whence $\Menv_\ell W(z_*)\le\Menv_\ell v(z_*)$ and, by
\noderef{EnvSlopeM} and $W(z_*)=v(z_*)$,
\begin{equation}\label{eq:ec-slope-minus}
  \Senv_\ell^-W(z_*)\ge\Senv_\ell^-v(z_*)\qquad(0<\ell<1).
\end{equation}

Fix $h\in(0,1)$ and $\varepsilon>0$. The second clause of
\noderef{EnvelopeSlopeOrder} applied to $W$ at $z_*$ with these $h$ and
$\varepsilon$ --- its hypothesis $1+\sigma>0$ for $|\sigma|\le h$ holds as
$h<1$ --- supplies $\delta>0$ with
$\Senv_\ell^-W(z_*)\le\Senv_h^+W(z_*)+\varepsilon$ for $0<\ell<\delta$. Pick
one $\ell$ with $0<\ell<\min\{\delta,1\}$. Then, using
\eqref{eq:ec-ham} and \eqref{eq:ec-slope-minus},
\[
  H=\inf_{\ell'>0}\Senv_{\ell'}^-v(z_*)\le\Senv_\ell^-v(z_*)
   \le\Senv_\ell^-W(z_*)\le\Senv_h^+W(z_*)+\varepsilon .
\]
As $\varepsilon>0$ was arbitrary, $H\le\Senv_h^+W(z_*)$ for every
$h\in(0,1)$; and the first clause of \noderef{EnvelopeSlopeOrder} applied to
$W$ --- whose hypothesis $1+\sigma>0$ for $\sigma\in[0,h']$ always holds ---
shows that $h\mapsto\Senv_h^+W(z_*)$ is nondecreasing on $(0,\infty)$, so
$H\le\Senv_h^+W(z_*)$ for every $h>0$. Conversely
\eqref{eq:ec-slope-plus} and \eqref{eq:ec-ham} give
$\inf_{h>0}\Senv_h^+W(z_*)\le\inf_{h>0}\Senv_h^+v(z_*)=H$. Hence
\[
  \inf_{h>0}\Senv_h^+W(z_*)=H,
\]
and since $h\mapsto\Senv_h^+W(z_*)$ is a nondecreasing real-valued function on
$(0,\infty)$ with infimum $H$, it converges to $H$ as $h\downarrow0$:
\[
  \lim_{\ell\downarrow0}\Senv_\ell^+W(z_*)=H=H_v(z_*).
\]

Collecting \eqref{eq:ec-kappa}, \eqref{eq:ec-final}, Step 6(a)--(e) and Step 7
gives all the assertions, with this $\kappa$, this $H$ and this $v$.
\end{proof}
